% ============================================================
% bidi-auto.tex  —  XeLaTeX template for plain-text mixed
% Latin / Hebrew input with ZERO markup in the source.
%
% The Unicode Bidirectional Algorithm (UBA) handles direction
% automatically. XeLaTeX + bidi respect it at the paragraph
% level. Font fallback via fontspec covers missing glyphs.
%
% Compile:
%   pandoc input.md -o output.pdf \
%     --pdf-engine=xelatex \
%     --template=bidi-auto.tex
%
% Or standalone:
%   xelatex bidi-auto.tex
% ============================================================

% \documentclass[12pt, a4paper]{article}

\usepackage{iftex}
\RequireXeTeX

% ------------------------------------------------------------
% SINGLE FONT covering Latin + Hebrew — the key to zero markup.
% Taamey Frank CLM or SBL BibLit both span the full range.
% fontspec will fall back to the FallbackFonts list for any
% glyph the main font lacks; order matters.
% ------------------------------------------------------------
\usepackage{fontspec}

% \setmainfont{FreeSans}[
%   Ligatures    = TeX,
%   FallbackFonts = {SBL BibLit, SBL Hebrew, FreeSans, DejaVu Sans}
% ]

% That's it for fonts. No \newfontfamily, no \texthebrew{}.

% ------------------------------------------------------------
% BIDI — loads the Unicode Bidirectional Algorithm support.
% With bidi loaded, XeLaTeX reads the Unicode script properties
% of each character and sets run direction automatically.
% No \LR or \RL commands needed in the source.
% ------------------------------------------------------------
\usepackage{fontspec}
% \defaultfontfeatures{Renderer=HarfBuzz}
\setmainfont{Source Sans 3}[
  Ligatures    = TeX,
  FallbackFonts = {Alef, Arial Hebrew, Lucida Grande, FreeSans, Source Sans 3}
]
% \usepackage{fontspec}
% \defaultfontfeatures{Renderer=HarfBuzz}

% \setmainfont{Avenir Next}

% \newfontfamily\hebrewfont{Arial Hebrew}[
%   Renderer = HarfBuzz,
%   Script   = Hebrew
% ]

\usepackage{fontspec}
\defaultfontfeatures{Renderer=HarfBuzz}

% \setmainfont{Avenir Next}

\newfontfamily\hebrewfont{Arial Hebrew}[
  Renderer = HarfBuzz,
  Script   = Hebrew
]

\usepackage[Hebrew]{ucharclasses}
% \setTransitionsFor{Hebrew}{\hebrewfont}{\normalfont}
% \usepackage[Hebrew, IPAExtensions]{ucharclasses}

\newfontfamily\ipafont{Source Sans 3}[Renderer=HarfBuzz]

\setTransitionsFor{Hebrew}{\hebrewfont}{\normalfont}
% \setTransitionsFor{IPAExtensions}{\ipafont}{\normalfont}
% \usepackage{ucharclasses}
% \setTransitionsForHebrew{\hebrewfont}{\normalfont}
\usepackage{bidi}
% \renewcommand{\,}{\thinspace}
\usepackage{etoolbox}
\AtBeginEnvironment{footnote}{\hbadness=10000\hfuzz=10pt}
\let\oldfootnote\footnote
% \renewcommand{\footnote}[1]{\oldfootnote{\LTR #1}}
\usepackage{etoolbox}
% \patchcmd{\footnote}{\ignorespaces}{\leavevmode\LTR\ignorespaces}{}{}
% \let\bidi@saved@everypar\everypar
\usepackage{endfloat}
\AtEndDocument{\let\LTR\relax\let\RTL\relax}
\AtBeginDocument{
  \let\LTR\relax
  \let\RTL\relax
}
\usepackage[hyperfootnotes=false]{hyperref}
